\section{Conclusion}
\label{sec:conclusion}

\begin{itemize}
  \item Calculus versus coherence space: they are equal at first order
  and differ from second order on. The difference is exactly the
  non-sequential behaviours.
  \item At second order the gap is governed by induced $4$-cycles, and it
  transfers along negative occurrences.
  \item Negation swaps the global and local calculi, and the coherence
  space is self-dual between them. The gap at $T$ and the Bell gap at the
  core are two sides of one gap.
  \item The certain-state calculus is not compositional (tight factors,
  non-tight tensor). The coherence space is compositional by
  construction, which is what it pays for with extra behaviours.
  \item Consequence: reasoning in the coherence space certifies weaker
  guarantees than sequential programs satisfy.
  \item Open questions: an exact condition beyond second order, a
  theta-body (quantum) calculus for programs, exponentials beyond
  $!\mathrm{Bool}$, a characterization of the tight formulas, whether
  $h$ is quantum (Ac\'in--Fritz--Leverrier--Sainz), and a comparison with
  Kissinger--Uijlen's causal processes.
\end{itemize}
